\documentclass[11pt]{article}
\usepackage[margin=1in]{geometry}
\usepackage{booktabs,graphicx,amsmath,amssymb}
\usepackage[hidelinks]{hyperref}
\renewcommand{\thetable}{S\arabic{table}}
\renewcommand{\thefigure}{S\arabic{figure}}

\title{Supplementary Material\\
Constraint-Driven Sparse Fragment Reconstruction for Digital Forensics}
\date{}

\begin{document}
\maketitle

This supplement holds evidence that supports the main paper without carrying a
headline claim. Every table is generated from the released CSVs by
\texttt{scripts/render\_stats\_tables.py} and
\texttt{scripts/update\_paper\_tables.py}, the same generators that populate the
main paper, so the numbers here cannot drift from the artefacts they come from.

\section{Constraint ablation}
The contribution of each constraint family, averaged over all 20 cells.


\begin{table*}[htbp]
\centering
\caption{Constraint ablation, averaged over the full 20-cell sweep and seeds $\{0,1,2\}$. Every variant retains the C1 data term, the $\ell_1$ sparsity prior, and the final C1 projection (so C1 $=0$ throughout, omitted). CVR columns in \%.}
\label{tab:ablation}
\setlength{\tabcolsep}{6pt}
\resizebox{0.95\textwidth}{!}{%
\begin{tabular}{lcccccc}
\toprule
\textbf{Variant} & \textbf{PSNR (dB)} & \textbf{SSIM} & \textbf{HRP} & \textbf{CVR C2} & \textbf{CVR C3} & \textbf{CVR C4} \\
\midrule
% BEGIN AUTO ABLATION
C1 only & 24.41 & 0.716 & 1.76 & 1.55 & 5.18 & 2.69 \\
C1+C2 (TV) & 25.56 & 0.752 & 1.76 & 1.59 & 0.08 & 2.08 \\
C1+C3 (bounds) & 25.09 & 0.727 & 1.58 & 1.52 & 0.00 & 2.66 \\
C1+C4 (freq.) & 24.07 & 0.725 & 1.83 & 1.68 & 5.38 & 1.69 \\
Full CSFR & 25.68 & 0.757 & 1.73 & 1.46 & 0.00 & 1.85 \\
% END AUTO ABLATION
\bottomrule
\end{tabular}%
}
\end{table*}

\section{Runtime and memory}
Wall-clock and peak memory per patch, measured on the CPU reference machine.


\begin{table}[htbp]
\centering
\caption{Per-patch runtime and peak memory, CPU, $64\times64$ patches. ms/patch and per-batch mean$\pm$SD over five repeats of a 32-patch batch; peak is \texttt{tracemalloc} peak allocation. Full environment in \texttt{results/bench/environment.json}.}
\label{tab:bench}
\setlength{\tabcolsep}{5pt}
\begin{tabular}{lccccc}
\toprule
\textbf{Method} & \textbf{ms/patch} & \textbf{mean (s)} & \textbf{SD (s)} & \textbf{median (s)} & \textbf{peak (MB)} \\
\midrule
% BEGIN AUTO BENCH
Zero-fill & 0.03 & 0.001 & 0.000 & 0.001 & 1.6 \\
Linear interpolation & 84.64 & 2.708 & 0.290 & 2.596 & 2.4 \\
Inpainting (Telea) & 1.61 & 0.052 & 0.006 & 0.049 & 1.5 \\
Inpainting (Navier--Stokes) & 1.18 & 0.038 & 0.001 & 0.038 & 1.5 \\
Dictionary learning & 1.46 & 0.047 & 0.001 & 0.047 & 1.5 \\
CSFR (proposed) & 44.35 & 1.419 & 0.105 & 1.435 & 1.0 \\
% END AUTO BENCH
\bottomrule
\end{tabular}
\end{table}

\section{Iteration-budget sensitivity}
CSFR PSNR against solver budget on a 100-patch, seed-0 subset. The gap column
compares against the strongest baseline's full-cell value, so it slightly
understates CSFR where the subset is easier than the full cell.


\begin{table}[htbp]
\centering
\caption{Iteration-budget sensitivity. CSFR PSNR (dB) is recomputed on a 100-patch, seed-0 subset at each budget; the ``Gap@4800'' column is CSFR at $T{=}4800$ minus the strongest baseline's full-cell (300-patch, seed-0) PSNR in that cell, so it slightly understates CSFR where the subset is easier than the full cell. Extending the budget closes most of the CF3/CF4 deficit and reverses the CF2 deficit, establishing that the reported 600-iteration losses on these families are substantially an under-convergence effect.}
\label{tab:budget}
\setlength{\tabcolsep}{5pt}
\begin{tabular}{lccccc}
\toprule
\textbf{Cell} & \textbf{T=600} & \textbf{T=1200} & \textbf{T=2400} & \textbf{T=4800} & \textbf{Gap@4800} \\
\midrule
% BEGIN AUTO BUDGET_SENS
CF1~L5 & 28.15 & 28.15 & 28.15 & 28.15 & +0.10 \\
CF2~L3 & 25.78 & 27.81 & 28.00 & 28.00 & +1.46 \\
CF2~L5 & 16.57 & 20.44 & 22.32 & 22.45 & +1.32 \\
CF3~L5 & 11.21 & 12.36 & 14.56 & 17.44 & -0.61 \\
CF4~L5 & 11.03 & 12.14 & 14.39 & 17.62 & -0.61 \\
% END AUTO BUDGET_SENS
\bottomrule
\end{tabular}
\end{table}

\section{Bounded-Unknown and the signal-covering subset}
Abstention rates and the continuous isolated-missing-fraction per cell. The
binary flag saturates by family and level, which is why the continuous
statistic is the one reported in the main paper.


\begin{table*}[htbp]
\centering
\caption{Bounded-Unknown and abstention rates by cell. The binary Bounded-Unknown flag at threshold $0.5$ on isolated-missing-fraction saturates: it is zero for every CF1 cell and for CF3/CF4~L1, unity for CF3/CF4 from L2 upward and CF2 from L3 upward, and varies within a cell only at CF2~L1 and L2. It is a family-and-level indicator rather than a per-patch signal. The continuous isolated-missing-fraction (columns: mean, SD, median) is reported beside it and should be preferred for analysis. The signal-covering subset count $n_{\text{sig}}$ is the number of patches in which at least a quarter of the centre box is erased, so that the corruption overlaps the digit; it is defined on the corruption geometry, not on constraint violations. It is empty for CF1 at L1--L2 and for CF3/CF4 at L1--L3, and saturates to 1000 for every family from L4 upward, leaving CF1~L3 and CF2~L1--L3 as the only cells with within-cell variation.}
\label{tab:downstream-abstain}
\setlength{\tabcolsep}{3pt}
\resizebox{0.95\textwidth}{!}{%
\begin{tabular}{llccccc}
\toprule
\textbf{Fam.} & \textbf{Lvl} & \textbf{BU rate} & \textbf{IMF mean} & \textbf{IMF SD} & \textbf{IMF median} & \textbf{$n_{\text{sig}}$} \\
\midrule
% BEGIN AUTO DOWNSTREAM_ABSTAIN
CF1 & L1 & 0.000 & 0.000 & 0.001 & 0.000 & 0 \\
CF1 & L2 & 0.000 & 0.001 & 0.003 & 0.000 & 0 \\
CF1 & L3 & 0.000 & 0.010 & 0.006 & 0.010 & 982 \\
CF1 & L4 & 0.000 & 0.070 & 0.012 & 0.070 & 1000 \\
CF1 & L5 & 0.000 & 0.253 & 0.016 & 0.252 & 1000 \\
CF2 & L1 & 0.880 & 0.612 & 0.091 & 0.619 & 352 \\
CF2 & L2 & 0.992 & 0.662 & 0.067 & 0.664 & 760 \\
CF2 & L3 & 1.000 & 0.710 & 0.050 & 0.713 & 991 \\
CF2 & L4 & 1.000 & 0.778 & 0.036 & 0.781 & 1000 \\
CF2 & L5 & 1.000 & 0.849 & 0.022 & 0.851 & 1000 \\
CF3 & L1 & 0.000 & 0.500 & 0.000 & 0.500 & 0 \\
CF3 & L2 & 1.000 & 0.800 & 0.000 & 0.800 & 0 \\
CF3 & L3 & 1.000 & 0.900 & 0.000 & 0.900 & 0 \\
CF3 & L4 & 1.000 & 0.938 & 0.000 & 0.938 & 1000 \\
CF3 & L5 & 1.000 & 0.955 & 0.000 & 0.955 & 1000 \\
CF4 & L1 & 0.000 & 0.500 & 0.000 & 0.500 & 0 \\
CF4 & L2 & 1.000 & 0.800 & 0.000 & 0.800 & 0 \\
CF4 & L3 & 1.000 & 0.900 & 0.000 & 0.900 & 0 \\
CF4 & L4 & 1.000 & 0.938 & 0.000 & 0.938 & 1000 \\
CF4 & L5 & 1.000 & 0.955 & 0.000 & 0.955 & 1000 \\
% END AUTO DOWNSTREAM_ABSTAIN
\bottomrule
\end{tabular}%
}
\end{table*}

\end{document}
